\section{ANEXOS}

Este anexo reúne lo que se necesita para auditar o reproducir el trabajo. La Tabla~\ref{tab:anexos} indica qué va en el documento y qué queda en el repositorio, con su ruta.

\begin{table}[htbp]
\centering
\caption{Índice de anexos y de los entregables que quedan en el repositorio}
\label{tab:anexos}
\small
\renewcommand{\arraystretch}{1.2}
\begin{tabular}{@{}p{4.0cm}p{9.0cm}@{}}
\toprule
\textbf{Material} & \textbf{Contenido y ruta en el repositorio} \\
\midrule
\multicolumn{2}{@{}l}{\emph{En este documento}}\\
Constantes, glosario & Tabla de constantes numéricas y glosario (a continuación) \\
Calibración y NL2C & Detalle de la calibración del optimizador y resumen de resultados del asistente NL2C \\
Datos, algoritmos, pantallas, ética & Datos del sistema, tres algoritmos del pipeline, cuatro pantallas y el protocolo ético con el texto de consentimiento del OE5 \\
\midrule
\multicolumn{2}{@{}l}{\emph{Documentos de la ingeniería (IEEE 830, 1016 y 1058, ATAM)}}\\
SRS, SDD, SPMP, ATAM & Especificación de requisitos, diseño, plan del proyecto e informe ATAM, en \path{docs/} \\
ADR & 49 registros de decisión: 0001 a 0049 en \path{docs/arquitectura/adr/} \\
Actas OE1 y OE2 & \path{docs/arquitectura/ACTA_APROBACION_OE1.md} y \path{ACTA_APROBACION_OE2.md}; sin firmas \\
Datos personales (Anexo M) & \path{docs/arquitectura/ANEXO_M_DATOS_PERSONALES.md}: tratamiento y base legal (Ley 1581 de 2012) \\
Antecedentes, requerimientos y MoSCoW & Conceptos de base, requerimientos con su justificación y priorización de los 98 requisitos, en \path{docs/requisitos/} y \path{docs/analisis/} \\
Autorización con ámbito & Auditoría de autorización y protección de datos, en \path{docs/arquitectura/} \\
Algoritmos restantes & Seis procedimientos del pipeline (percepción, normalización, puntaje, resolución, frente, selección), en \path{docs/algoritmos/} \\
Galería de pantallas & Las 29 capturas de la versión integrada, en \path{docs/entrega/} \\
Ampliaciones técnicas & Ampliaciones de los capítulos V y VII, en \path{docs/informes/} \\
\midrule
\multicolumn{2}{@{}l}{\emph{Evidencia}}\\
Hallazgos de calidad & 85 hallazgos vigentes (16 de severidad alta, 58 media y 11 baja, según el registro), en \path{docs/informes/calidad/hallazgos_qa_oficial.csv} \\
Instrumentos del OE5 & \path{docs/validacion/oe5/anexos_OE5/} y \path{docs/validacion/oe5/oe5_cierre/formulario_cierre.md} \\
Modelo de negocio & \path{docs/negocio/punto_equilibrio.py} \\
\bottomrule
\end{tabular}
\end{table}

Los diagramas de contexto y de contenedores de la alternativa elegida son las
Figuras~\ref{fig:c4-contexto} y~\ref{fig:c4-contenedores}, y el protocolo ético y los
instrumentos corresponden al OE5 (Sección~\ref{sec:oe5}).

\subsection{Tabla de constantes numéricas}
\label{anx:constantes}

La Tabla~\ref{tab:constantes} reúne cada constante numérica del sistema con su valor, su módulo, su origen y su estado, según se explica en la Sección~\ref{sec:constantes}.

{\footnotesize
\setlength{\tabcolsep}{3pt}
\renewcommand{\arraystretch}{1.1}
\begin{longtable}{@{}p{2.6cm}p{4.3cm}p{1.9cm}p{3.7cm}p{1.5cm}@{}}
\caption{Constantes numéricas del sistema, con su origen y su estado.}\label{tab:constantes}\\
\toprule
Constante & Valor & Módulo & Origen & Estado \\
\midrule
\endfirsthead
\toprule
Constante & Valor & Módulo & Origen & Estado \\
\midrule
\endhead
\bottomrule
\endfoot
\multicolumn{5}{@{}l}{\emph{Percepción}}\\
Tope de páginas & 40 & percepción & Experimento: el PDF más largo del corpus tiene 39 páginas & Vigente \\
Temperatura del modelo & 0 & percepción & Decodificación voraz, para reducir la variabilidad entre lecturas & Vigente \\
Tiempo límite por llamada & 40 s & percepción & Una página típica tarda unos 5 s; 40 s indica una llamada colgada & Vigente \\
Reintentos y espera & 1 reintento, 3 s & percepción & Solo ante fallas de transporte, para no ocultar respuestas mal formadas & Vigente \\
\midrule
\multicolumn{5}{@{}l}{\emph{Normalización y resolución}}\\
Dimensión del vector & 1.024 & catálogo & Fija la columna de la base; se pide explícitamente al modelo de \emph{embeddings} & Vigente \\
Lote de \emph{embeddings} & 100 textos & catálogo & Una petición por lote, para evitar cuerpos de petición excesivos & Vigente \\
Umbral de reconocimiento & 0,5 & catálogo & Corte por debajo del cual el renglón no corresponde a ningún concepto; fijado al corregir el caso «chaquiras» & Vigente \\
Aporte del coseno $\gamma$ & 0,15 & catálogo & Ajuste sobre casos de error; pequeño para que el coseno desempate y no decida & En ciclo de calibración \\
Umbral de rescate $\theta$ & 0,86 & catálogo & Ajuste sobre casos de error; separa un sinónimo de un parecido de tema & En ciclo de calibración \\
Piso del rescate $\phi$ & 0,55 & catálogo & Apenas sobre el corte de 0,5, para que un rescate entre a revisión & En ciclo de calibración \\
Umbral de autoaceptación $\tau$ & 0,75 & inteligencia & Valor por omisión (\texttt{TUKIT\_TAU}); sin decisiones humanas para calibrarlo, ver Sección~\ref{sec:optimizacion-cal} & Sin calibrar, declarado \\
Tolerancia de medidas & 0,25 y rango 0,85 a 1,35 & catálogo & Ajuste con las auditorías de los kits (Tabla~\ref{tab:kits}) & Vigente \\
Rango de precio en tiendas & \$1.000 a \$400.000 & catálogo & Descarta accesorios sueltos y lotes institucionales & Vigente \\
\midrule
\multicolumn{5}{@{}l}{\emph{Optimización}}\\
Límite de enumeración & 4.096 ($2^{12}$) & inteligencia & Medido: la enumeración alcanza a NSGA-II en $S\approx4.932$ (mediana) y $2.380$ (percentil 95); ver Tabla~\ref{tab:calopt-umbral} & Vigente \\
Población & $\min(100,\max(20,6n))$, par & inteligencia & Mitad de la fórmula de partida, elegida por validación cruzada agrupada por colegio (Tabla~\ref{tab:calopt-rejilla}); el par evita desperdiciar un individuo & Vigente \\
Generaciones &\raggedright $\min(60,\max(15,\,10\log_{10}|\mathcal{E}|))$ & inteligencia & Escala con el logaritmo del espacio $|\mathcal{E}|$; mitad de la de partida, misma calibración & Vigente \\
Cruce SBX & $p_c=0{,}9$, $\eta_c=15$ & inteligencia & Valores estándar de Deb y Agrawal \citep{deb_simulated_1995} & Vigente \\
Mutación polinomial & $p=1$, $\eta_m=20$ & inteligencia & Valores estándar de Deb y Goyal \citep{deb_combined_1996} & Vigente \\
Semilla & 0 & inteligencia & Reproducibilidad del frente entre ejecuciones & Vigente \\
Pesos por omisión (costo, marca, calidad, personaje) & 0,285; 0,273; 0,255; 0,187 & inteligencia & Encuesta, pregunta P21 (n=60), con IC 95\,\% \emph{bootstrap}; reemplazan 0,35; 0,30; 0,25; 0,10 (Sección~\ref{sec:optimizacion-cal}) & Vigente \\
Sobreprecio para Premium & 10\,\% & inteligencia & Encuesta: 42 de 60 pagan como máximo un 5\,\% más; con al menos 10\,\%, MAUT decide si se marca el Premium & Vigente \\
Opciones por renglón & 3 & catálogo & Pocas opciones diferenciadas facilitan la decisión \citep{iyengar_when_2000} & Vigente \\
\midrule
\multicolumn{5}{@{}l}{\emph{Requisitos de desempeño y operación}}\\
Tiempo de generación del kit & 8 s & SRS (DP-03) & Requisito de desempeño; se contrasta en la Sección~\ref{sec:optimizacion} & Vigente \\
Ventana de revisión & 4 h en temporada & SRS (ML-09) & Requisito: el administrador revisa cada lista nueva antes de publicarla & Vigente \\
Disyuntor & 5 fallas, 30 s & biblioteca común & Patrón de resiliencia; 30 s es el escenario ESC-R2 del ATAM & Vigente \\
Caché de lecturas & 30 días & inteligencia & Evita pagar de nuevo la percepción de un mismo documento & Vigente \\
Caché de vectores & 24 h & catálogo & Evita recalcular el vector de un renglón repetido & Vigente \\
\end{longtable}
}

\subsection{Glosario}
\label{anx:glosario}

\begin{description}[leftmargin=2.5cm,style=nextline]
\item[MAUT] Teoría de utilidad multiatributo: método que resume varios criterios en un solo valor de utilidad mediante pesos que suman uno.
\item[NSGA-II] Algoritmo genético de ordenamiento no dominado que aproxima el frente de Pareto de un problema multiobjetivo \citep{deb_fast_2002}.
\item[Frente de Pareto] Conjunto de soluciones no dominadas, es decir, aquellas en las que ninguna otra mejora un objetivo sin empeorar otro.
\item[Hipervolumen] Volumen del espacio de objetivos que un frente domina respecto de un punto de referencia; mide la calidad de un conjunto de soluciones \citep{zitzler_multiobjective_1999}.
\item[NL2C] Traducción de una preferencia en lenguaje natural a restricciones estructuradas que el optimizador entiende.
\item[$\tau$] Umbral de autoaceptación: puntaje por encima del cual un renglón se da por resuelto sin revisión individual.
\item[$\delta$] Margen mínimo entre el mejor y el segundo candidato que usó una versión anterior del umbral.
\item[$\kappa$] Coeficiente kappa de Cohen: acuerdo entre dos clasificadores corregido por el azar \citep{cohen_coefficient_1960}.
\item[SUS] Escala de usabilidad del sistema, de 0 a 100 \citep{brooke_sus_1996}.
\item[NPS] Índice de recomendación neta, de $-100$ a $100$ \citep{reichheld_one_2003}.
\item[RTLX] Versión sin ponderar del índice de carga de trabajo NASA-TLX, de 0 a 100 \citep{hart_development_1988,hart_nasa_2006}.
\item[CRISP-DM] Proceso de seis etapas para proyectos de datos \citep{chapman_crisp-dm_2000,wirth_crisp-dm_2000}.
\item[ATAM] Método de evaluación de arquitecturas por compromisos entre atributos de calidad \citep{bass_software_2021}.
\item[ADR] Registro de decisión de arquitectura.
\item[Concepto canónico] Clase de productos equivalentes del catálogo a la que se asigna un renglón de la lista.
\end{description}
\subsection{Tablas de apoyo de los capítulos}
\label{anx:tablas}

En este anexo se reúnen las tablas extensas a las que remiten los capítulos III, V y VII. Cada capítulo explica en su texto la conclusión que la tabla respalda.

\begin{table}[htbp]
\centering
\caption{Reparto de tablas por contexto delimitado}
\label{tab:reparto}
\small
\begin{tabular}{@{}lrp{8.4cm}@{}}
\toprule
\textbf{Servicio} & \textbf{Tablas} & \textbf{Responsabilidad} \\
\midrule
Identidad     & 16 & Personas, tenencia, autorización y consentimiento. Frontera de cumplimiento \\
Catálogo      & 19 & Qué productos existen y a qué precio público. Reconstruible desde sus archivos de datos (CF-02a) \\
Inteligencia  & 12 & Percepción, normalización, resolución y optimización \\
Comercio      & 18 & Pedidos, inventario, pagos y entrega. Irremplazable \\
Operación     & 9  & Casos, notificaciones, auditoría y telemetría. Solo anexar \\
\midrule
\textbf{Total} & \textbf{74} & \\
\bottomrule
\end{tabular}
\end{table}

\begin{table}[htbp]
\centering
\caption{Políticas inejecutables y su corrección en el esquema}
\label{tab:politicas}
\small
\begin{tabular}{@{}p{4.3cm}p{5.0cm}p{4.4cm}@{}}
\toprule
\textbf{Política enunciada} & \textbf{Por qué era inejecutable} & \textbf{Corrección} \\
\midrule
Revocar el consentimiento del tratamiento de un menor &
La tabla registraba que el acudiente autorizó el tratamiento de un hijo, pero ninguna
columna indicaba de cuál. Con dos hijos, una fila los cubría de forma indistinguible &
Referencia a la matrícula, con restricción de coherencia contra el régimen \\
\addlinespace
Destruir la llave de la matrícula 12 meses después del cierre &
No existía columna que registrara el cierre &
Fecha de cierre de la matrícula \\
\addlinespace
Informar al titular que su dato fue suprimido &
El proceso marcaba la solicitud como completada tras anonimizar la fila local, sin esperar
confirmación de ningún otro servicio &
Estado de propagación y registro de los servicios que confirmaron \\
\addlinespace
Conservar evidencia de lo revocado &
Una actualización sobrescribía la fila, eliminando la constancia de qué se autorizó y
cuándo &
Columnas de vigencia con índice único parcial sobre la fila vigente \\
\bottomrule
\end{tabular}
\end{table}

\begin{table}[htbp]
\centering
\caption{Desglose de latencia por etapa (lista de 20 ítems, mediana)}
\label{tab:latencia}
\small
\begin{tabular}{@{}p{6.2cm}rrp{3.2cm}@{}}
\toprule
\textbf{Etapa} & \textbf{p50 (ms)} & \textbf{\% total} & \textbf{Naturaleza} \\
\midrule
Expansión por grafo de sustituibilidad & 33.470,0 & 95,28 & CPU en Python \\
Embedding de la consulta & 832,0 & 2,37 & CPU nativo (torch) \\
Optimización NSGA-II & 362,2 & 1,03 & CPU en Python \\
Percepción por reconocimiento óptico & 183,2 & 0,52 & CPU nativo \\
Consultas a la base de datos & 274,3 & 0,78 & Espera de E/S \\
Resto & 5,5 & 0,02 & \\
\midrule
\textbf{Total} & \textbf{35.127,2} & \textbf{100,00} & \\
\bottomrule
\end{tabular}
\end{table}

\begin{table}[htbp]
\centering
\caption{Registros de decisión arquitectónica estructurales}
\label{tab:adr}
\small
\begin{tabular}{@{}cp{5.0cm}p{7.6cm}@{}}
\toprule
\textbf{ADR} & \textbf{Decisión} & \textbf{Fundamento} \\
\midrule
0011 & Separar el catálogo del comercio & Elimina el doble maestro del catálogo. Los precios
no se unifican: el costo contractual y el precio público observado son hechos distintos en
lados opuestos de la asimetría de recuperación \\
\addlinespace
0012 & Matrícula como ancla de tenencia y roles con ámbito & Un acudiente puede tener hijos
en colegios distintos: la tenencia no reside en la persona \\
\addlinespace
0013 & Monorrepositorio con despliegue independiente & Separación entre vista de
desarrollo y vista física; fronteras instrumentadas con funciones de aptitud \\
\addlinespace
0014 & Dirección de entrega: libreta y copia congelada & Un hecho histórico y un dato
editable son entidades distintas aunque se escriban igual \\
\addlinespace
0015 & Dónde se evalúa el ámbito cuando el ancla vive en otra base & Una política de fila
no puede unirse contra otra base: se copia la clave de ámbito, no la lista de matrículas, y
se deniega por omisión. Sigue siendo diseño (Sección~\ref{sec:defectos}) \\
\addlinespace
0037 & Eventos y trabajos por RabbitMQ, con la base de datos como respaldo & Una cola por
consumidor, con reintentos y descartados; sin RabbitMQ, el sistema sigue por sondeo de las
bandejas de salida \\
\bottomrule
\end{tabular}
\end{table}

\begin{table}[htbp]
\centering\footnotesize
\setlength{\tabcolsep}{2pt}
\renewcommand{\arraystretch}{1.35}
\newcommand{\enc}[2]{{\scriptsize\begin{tabular}[b]{@{}c@{}}#1\\#2\end{tabular}}}
\caption{Comparación de Tukit con las soluciones existentes.}
\label{tab:comparacion-soluciones}
\begin{tabular}{@{}>{\raggedright\arraybackslash}p{3.4cm}
  >{\centering\arraybackslash}p{1.25cm}
  *{5}{>{\centering\arraybackslash}p{1.25cm}}
  >{\columncolor{gray!15}\centering\arraybackslash}p{1.25cm}@{}}
\toprule
\enc{}{\textbf{Capacidad}} &
\enc{Teacher-}{Lists} &
\enc{1st Day}{Supplies} &
\enc{Pana-}{mericana} &
\enc{Grupo}{Éxito} &
\enc{Papelería}{Primavera} &
\enc{Mercado}{Libre} &
\enc{\textbf{Tukit}}{} \\
\midrule
Parte de la lista oficial del colegio
  & \cumple & \cumple & \nocumple & \nocumple & \nocumple & \nocumple & \cumple \\
Lee listas en foto, PDF o texto
  & \nocumple & \nocumple & \nocumple & \nocumple & \nocumple & \nocumple & \cumple \\
Reconoce el vocabulario escolar local
  & \nocumple & \nocumple & \nocumple & \nocumple & \nocumple & \nocumple & \cumple \\
Arma el kit completo
  & \parcial & \cumple & \nocumple & \parcial & \parcial & \parcial & \cumple \\
Ofrece un catálogo general, más allá de la lista escolar
  & \parcial & \nocumple & \cumple & \cumple & \cumple & \cumple & \nocumple \\
Considera productos de varias tiendas
  & \parcial & \nocumple & \nocumple & \nocumple & \nocumple & \parcial & \cumple \\
Controla directamente su inventario y disponibilidad
  & \nocumple & \cumple & \cumple & \cumple & \cumple & \parcial & \nocumple \\
Ofrece varias alternativas según presupuesto y preferencias
  & \nocumple & \nocumple & \nocumple & \parcial & \nocumple & \parcial & \cumple \\
Disponible en Colombia
  & \nocumple & \nocumple & \cumple & \cumple & \cumple & \cumple & \cumple \\
\bottomrule
\end{tabular}

\vspace{2pt}
{\footnotesize \cumple\ Cumple \quad \parcial\ Cumple parcialmente \quad \nocumple\ No cumple}
\end{table}

\begin{table}[htbp]
\centering
\caption[Auditorías de los kits del Colegio Agustiniano Ciudad Salitre]{Auditorías de los kits del Colegio Agustiniano Ciudad Salitre, 11 y 12 de octubre de
2026; las dos últimas, posteriores a la versión evaluada (cifras reportadas; costo de los trece kits de cada perfil)}
\label{tab:kits}
\small
\begin{tabular}{@{}lrrrrr@{}}
\toprule
\textbf{Auditoría} & \textbf{Comprables} & \textbf{Huecos} & \textbf{Cobros} &
\multicolumn{2}{c}{\textbf{Costo (millones de pesos)}} \\
\cmidrule(l){5-6}
 & \textbf{que cumplen} & \textbf{declarados} & \textbf{sin cumplir} & \textbf{Ahorrador} &
\textbf{Premium} \\
\midrule
1 (11 oct.) & 263 de 888 (29,6\,\%) & 240 & 589 & 13,9 & 32,0 \\
2 (11 oct.) & 465 de 888 (52,4\,\%) & 385 & 242 & 9,9 & 17,6 \\
3 (11 oct.) & 617 de 886 (69,6\,\%) & 220 & 49 & 9,0 & 14,2 \\
4 (12 oct.) & 758 de 885 (85,6\,\%) & 69 & 58 & 8,1 & 17,3 \\
\bottomrule
\end{tabular}
\end{table}

\begin{table}[htbp]
\centering
\caption{Impulsores arquitectónicos de la descomposición en servicios}
\label{tab:impulsores}
\small
\begin{tabular}{@{}p{3.1cm}p{5.0cm}p{5.6cm}@{}}
\toprule
\textbf{Criterio} & \textbf{Fundamento} & \textbf{Qué separa} \\
\midrule
Recuperabilidad asimétrica &
El objetivo de punto de recuperación difiere en órdenes de magnitud entre almacenes &
El catálogo se reconstruye desde sus archivos de datos, una vez exportada al repositorio la
curaduría del panel, y admite perder un día de cambios (CF-02a); un pedido no puede perder un
solo byte. Compartir almacén obliga a dimensionar el
respaldo por el más exigente y a exponer el crítico a la carga de escritura del volátil \\
\addlinespace
Frontera de cumplimiento &
Ley 1581 de 2012 \citep{congreso_ley1581_2012} y Decreto 1377 de 2013
\citep{presidencia_decreto1377_2013} &
Los datos de menores de edad se aíslan en un servicio con base, proceso y llaves propios.
La superficie que un auditor debe examinar se reduce a un solo componente \\
\addlinespace
Perfil de recurso &
Heterogeneidad de la demanda computacional &
La inferencia de modelos y la optimización multiobjetivo son intensivas en CPU y memoria;
la gestión transaccional requiere disponibilidad y consistencia fuerte. Colocarlas juntas
obliga a dimensionar todo por la más exigente \\
\bottomrule
\end{tabular}
\end{table}

\begin{table}[htbp]
\centering
\caption{Alternativas arquitectónicas evaluadas}
\label{tab:alternativas}
\small
\begin{tabular}{@{}p{2.6cm}p{3.5cm}p{3.5cm}p{4.1cm}@{}}
\toprule
& \textbf{A. Monolito modular} & \textbf{B. Microservicios de grano fino} & \textbf{B$'$. Grano grueso (elegida)} \\
\midrule
Componentes & 1 desplegable, módulos internos & 13 servicios y puerta de enlace & 5 servicios de dominio y puerta de enlace, 5 bases \\
\addlinespace
Aislamiento de datos de menores & Por convención de código & Por servicio & Por servicio, base y llaves \\
\addlinespace
Punto de recuperación & Único para todo & Por servicio & Diferenciado por criticidad \\
\addlinespace
Complejidad operativa & Baja & Alta: coordinación distribuida en cada transacción & Media \\
\addlinespace
Costo de infraestructura & Mínimo & Alto & Moderado \\
\addlinespace
Puntaje ponderado (sobre 5,00) & 2,82 & 3,44 & 3,91 \\
\bottomrule
\end{tabular}
\end{table}

\begin{table}[H]
\centering\footnotesize
\setlength{\tabcolsep}{4pt}
\renewcommand{\arraystretch}{1.1}
\caption{Entregables, objetivos, estándares asociados y justificación.}
\label{tab:entregables}
\begin{tabularx}{\textwidth}{@{}p{3cm}p{1.4cm}p{3.2cm}X@{}}
\toprule
Entregable & Objetivo & Estándares asociados & Justificación \\
\midrule
Memoria de trabajo de grado &  &
Formato de trabajo de grado del programa de Ingeniería de Sistemas &
Reúne en un solo documento el problema, la solución, su desarrollo y los resultados de los
cinco objetivos específicos. \\
Dossier de negocio & OE1 &
Design Thinking \citep{brown_design_2008}, lienzo del modelo de negocio
\citep{osterwalder_business_2010} &
Documenta el entendimiento de las familias y de los colegios, con las entrevistas, los mapas
de empatía, la propuesta de valor, el modelo de negocio y el análisis FODA. \\
Especificación de requerimientos (SRS) & OE1 &
IEEE 830 \citep{ieee_830_1998}, ISO/IEC/IEEE 29148 \citep{iso_29148_2018} &
Define qué debe hacer el sistema, con requerimientos funcionales y no funcionales
priorizados con el método MoSCoW y relacionados con las restricciones legales y de
negocio. \\
Documento de Diseño Maestro (ATAM, SDD, SAD) & OE2 &
Richards y Ford \citep{richards_fundamentals_2025}, ISO/IEC 25010 \citep{iso_25010_2011},
ATAM \citep{bass_software_2021}, IEEE 1016 \citep{ieee_1016_2009}, ISO/IEC/IEEE 42010
\citep{iso_42010_2022}, vistas 4+1 \citep{kruchten_41_1995} &
Documenta la arquitectura por vistas, los atributos de calidad que guiaron su diseño y la
comparación entre las tres alternativas evaluadas, de modo que cada decisión tenga su
razón explícita. \\
Anexo de experimentación & OE3, OE4 y OE5 &
CRISP-DM \citep{chapman_crisp-dm_2000} &
Registra, etapa por etapa, los experimentos de los componentes que dependen de datos, con
los datos usados, las métricas, los resultados y las decisiones que se tomaron a partir de
ellos. \\
Sistema Tukit & &
ISO/IEC 25010 \citep{iso_25010_2011} &
Implementa la solución en seis servicios con sus propias bases de datos, una aplicación
web para el acudiente y una aplicación de administración para el equipo de Tukit. \\
Documento de pruebas & OE5 &
ISO/IEC 25010 \citep{iso_25010_2011} &
Registra las pruebas que verifican cada componente y su integración en cada cambio del
código, de modo que un error se detecte antes de llegar a la versión entregada. \\
Informe de evaluación con usuarios y consentimiento informado & OE5 &
Lazar et al. \citep{lazar_research_2017}, SUS \citep{brooke_sus_1996}, NPS
\citep{reichheld_one_2003}, Ley 1581 de 2012 \citep{congreso_ley1581_2012}, Decreto 1377
de 2013 \citep{presidencia_decreto1377_2013} &
Documenta la evaluación de la usabilidad y de la carga cognitiva con acudientes, y el
consentimiento que exige la ley colombiana para tratar datos personales, incluidos los de
menores de edad. \\
\bottomrule
\end{tabularx}
\end{table}

\begin{table}[H]
\centering\footnotesize
\setlength{\tabcolsep}{4pt}
\renewcommand{\arraystretch}{1.15}
\caption{Distribución de los requerimientos por componente y prioridad.}
\label{tab:moscow}
\begin{tabular}{@{}lccclccc@{}}
\toprule
\textbf{Componente} & \textbf{M} & \textbf{S} & \textbf{C} &
\textbf{Componente} & \textbf{M} & \textbf{S} & \textbf{C} \\
\midrule
Catálogo & 5 & 1 & 0 & Pedidos y pagos & 10 & 3 & 0 \\
Percepción & 7 & 0 & 0 & Desempeño & 2 & 3 & 0 \\
Normalización & 7 & 1 & 0 & Restricciones de diseño & 1 & 0 & 0 \\
Optimización & 5 & 3 & 0 & Atributos de calidad & 18 & 1 & 1 \\
Aplicación del acudiente & 6 & 1 & 0 & Modelo de datos & 12 & 1 & 0 \\
Panel de administración & 6 & 1 & 0 & Despacho & 0 & 1 & 0 \\
Listas y matrículas & 2 & 0 & 0 & \textbf{Total (98)} & \textbf{81} & \textbf{16} &
\textbf{1} \\
\bottomrule
\end{tabular}
\end{table}


% Fuera del PDF (contenido en el repositorio, ver tab:anexos): anexo_figuras (copias en el cuerpo), anexo_antecedentes,
% anexo_requerimientos, anexo_autorizacion, anexo_hallazgos_tabla (docs/informes/calidad/hallazgos_qa_oficial.csv),
% anexo_ampliaciones (docs/informes/).
\subsection{Protocolo ético y consentimiento informado (OE5)}
\label{anx:etica}
Las sesiones del OE5 (Sección~\ref{sec:oe5}) se hicieron con madres, padres y acudientes mayores de 18 años del Colegio Agustiniano Ciudad Salitre. No participaron menores, y no se recogió ningún dato del hijo salvo el nivel y el grado. Por eso no se requirió asentimiento del menor; el artículo 7 de la Ley 1581 de 2012 se atiende evitando recoger datos de menores. Se excluyó a familiares del equipo y a quien tuviera relación laboral con el colegio o con el equipo. La participación fue voluntaria y sin pago, y el colegio solo transmitió la invitación sin recibir lista de participantes ni resultados individuales. El marco normativo es la Ley 1581 de 2012 (artículos 9 y 12) y el Decreto 1377 de 2013; el tratamiento de datos se detalla en el Anexo M (Tabla~\ref{tab:anexos}). El protocolo completo está en \path{docs/validacion/oe5/anexos_OE5/protocolo_etico.md}.

\paragraph{Datos y retención.} Las respuestas se guardan con un código de participante y no con el nombre; el nombre y el consentimiento se archivan aparte. No se grabó video ni se tomaron fotografías, y la pantalla y el audio solo se grabaron con autorización separada. La sesión corrió con una cuenta de prueba y pago simulado, sin tarjeta real. Las respuestas se conservan hasta la aprobación del trabajo de grado y 12 meses más, y luego se borran. Las grabaciones se borran al verificar el análisis y, a más tardar, a los 30 días. Los resultados se publican en conjunto y sin nombres. Si ocurre una filtración, el equipo informa a los afectados y a la dirección en un máximo de 72 horas.

\paragraph{Texto de consentimiento entregado al acudiente (versión 2, 2 de noviembre de 2026).}
Se envió por mensaje antes de la llamada, y el acudiente respondió con una plantilla que marca cada parte por separado, o lo firmó. A continuación se reproduce de forma condensada.

\begin{quote}\small
\textbf{TUKIT: prueba con acudientes.} Somos estudiantes de Ingeniería de Sistemas de la Pontificia Universidad Javeriana y hacemos nuestro trabajo de grado sobre TUKIT, una aplicación que arma la lista de útiles escolares y la convierte en un kit listo para comprar. Le pedimos unos 15 minutos para probarla con nosotros por videollamada.

\emph{Qué haría usted.} Desde su celular entra a la aplicación con una cuenta de prueba, hace tareas cortas (elegir colegio y grado, contar qué le gusta a su hijo, elegir un kit, cambiar un producto, pagar y ver el pedido) y dice en voz alta lo que piensa. Al final llena un formulario de unos 5 minutos. Se evalúa la aplicación, no a usted: no hay respuestas buenas ni malas.

\emph{Qué no pasa.} No se piden datos de su hijo o hija y él o ella no participa. El pago es simulado: no use su tarjeta real y no se le cobra nada. No se graba su cara ni su cámara.

\emph{Riesgos y beneficios.} El riesgo es mínimo: el tiempo de la llamada y alguna pregunta incómoda, que puede no responder. No hay pago. Su opinión ayuda a mejorar la aplicación para otros acudientes.

\emph{Voluntariedad.} Puede decir que no, saltarse una pregunta o parar en cualquier momento, sin explicaciones y sin consecuencia alguna.

\emph{Sus datos.} Se usan solo para el trabajo de grado y para mejorar el prototipo, y no se comparten con el colegio, con tiendas ni con terceros. Se guardan con un código de participante, y los resultados se publican en conjunto. Usted puede conocer, actualizar y corregir sus datos, pedir que se borren y retirar la autorización cuando quiera, escribiendo a \texttt{tukit.ayuda@gmail.com}.

\emph{Responsables.} El equipo de estudiantes de Ingeniería de Sistemas de la Pontificia Universidad Javeriana (Bogotá), bajo la dirección del director del trabajo de grado. Proyecto ISIS2610-G19.

\emph{Autorizaciones.} \textbf{Parte A (obligatoria para participar):} leí este texto, pude hacer preguntas, soy mayor de 18 años y autorizo de manera libre, previa e informada el uso de mis respuestas con los fines descritos. \textbf{Parte B (opcional):} autorizo grabar la pantalla del celular durante las tareas, solo para contar pasos y tiempos. \textbf{Parte C (opcional):} autorizo grabar el audio de mi voz, solo para transcribir lo que diga durante las tareas. Si responde que no a B o a C, igual puede participar y solo se toman notas.
\end{quote}

El formulario de cierre repite la Parte A al final como segundo registro, que no reemplaza al anterior a la llamada. La hoja de firma tiene una casilla de sí y no por cada parte, el nombre y la firma del acudiente, el código de participante y la firma de quien recibe el consentimiento.

\subsection{Pantallas del sistema construido}
\label{anx:pantallas}

Se muestran cuatro pantallas de las dos vistas: la espera y el pago del acudiente, y el seguimiento y la revisión de una lista.
Las capturas son de la versión integrada del 3 de noviembre de 2026, con datos de demostración; las
direcciones se reemplazaron por una genérica y el mapa se pixeló en las capturas 07a (dirección de entrega) y 13a (pedidos y mapa de entregas).

\newcommand{\pantalla}[3]{%
\begin{figure}[H]\centering\includegraphics[#1]{desarrollo/#2}\caption{#3}\end{figure}}
\newcommand{\pantallados}[5]{%
\begin{figure}[H]\centering
\begin{minipage}[t]{0.49\textwidth}\centering\includegraphics[width=\linewidth,height=0.36\textheight,keepaspectratio]{desarrollo/#1}\\\footnotesize (a) #2\end{minipage}\hfill
\begin{minipage}[t]{0.49\textwidth}\centering\includegraphics[width=\linewidth,height=0.36\textheight,keepaspectratio]{desarrollo/#3}\\\footnotesize (b) #4\end{minipage}
\caption{#5}\end{figure}}

\pantallados{04-avance-procesamiento-etapas}{Etapas del procesamiento}{04b-espera-revision-humana}{La
lista espera a una persona}{Pantallas de espera de una solicitud: lectura, normalización y búsqueda
terminadas con el armado por comenzar (a), y el aviso «Lo está revisando una persona» cuando la lista
pasa por la cola (b).}
\pantallados{07c-pasarela-de-pago}{Pasarela}{07d-pago-confirmado}{Pago confirmado}{Pantallas del pago en
modo simulado: la apertura de la pasarela (a) y el pedido pagado con su resumen (b).}
\pantallados{08-seguimiento-pedido}{Seguimiento}{08b-derecho-retracto}{Plazo de retracto}{Pantallas del
pedido: el seguimiento hasta la entrega, con la hora de cada paso (a), y el plazo del derecho de
retracto en el detalle del pedido (b).}
\pantallados{11a-revision-pdf-y-lectura}{PDF y lectura del modelo}{11b-renglones-con-concepto}{Renglones
con su concepto}{Pantallas de la revisión de una lista de Cuarto: el PDF original junto a lo que leyó
el modelo (a) y los renglones con su concepto y los atributos que exigen, entre ellos «un compás»
enviado a revisión (b).}

\iffalse % Galería completa (29 capturas): docs/entrega/ y carpeta figuras/desarrollo/

\pantalla{width=0.72\textwidth}{01-registro-acudiente}{Pantalla de registro en el modo «Soy nuevo
usuario», con los datos del acudiente y la casilla de la política de datos y los términos, que
enlaza a las páginas legales.}
\pantallados{02-hijo-colegio-grado}{Colegio y grado del hijo}{03-lista-oficial-o-subir}{Grado sin
lista verificada}{Pantallas de la matrícula del hijo: el colegio aliado y el grado (a) y, cuando el
grado no tiene lista verificada, la opción de subirla (b).}
\pantallados{04-avance-procesamiento-etapas}{Etapas del procesamiento}{04b-espera-revision-humana}{La
lista espera a una persona}{Pantallas de espera de una solicitud: lectura, normalización y búsqueda
terminadas con el armado por comenzar (a), y el aviso «Lo está revisando una persona» cuando la lista
pasa por la cola (b).}
\pantalla{width=0.72\textwidth,trim=0 464 0 0,clip}{05-kits-propuestos}{Pantalla de los kits del hogar
con el total de cada uno, el Fan recomendado con su razón, los productos con el personaje y el botón
del asistente del kit.}
\pantallados{06a-cambiar-producto-opciones}{Opciones del renglón}{15-asistente-propuesta}{Propuesta
del asistente}{Pantallas de ajuste del kit: el carrusel de opciones del renglón «Biblia católica», con
la diferencia de cada una en el total y el botón para descartar la línea (a), y el asistente del kit
con una propuesta que cambia un cuaderno de personaje por uno Norma y ahorra \$11.149 (b).}
\pantallados{07a-checkout-entrega}{Dirección confirmada en el mapa}{07b-checkout-formulario-lleno}{Datos
de entrega}{Pantallas de la entrega: la dirección buscada en OpenStreetMap con el punto confirmado en
el mapa (a) y el formulario completo antes de pagar (b).}
\pantallados{07c-pasarela-de-pago}{Pasarela}{07d-pago-confirmado}{Pago confirmado}{Pantallas del pago en
modo simulado: la apertura de la pasarela (a) y el pedido pagado con su resumen (b).}
\pantallados{08-seguimiento-pedido}{Seguimiento}{08b-derecho-retracto}{Plazo de retracto}{Pantallas del
pedido: el seguimiento hasta la entrega, con la hora de cada paso (a), y el plazo del derecho de
retracto en el detalle del pedido (b).}
\pantallados{09a-pqrs-formulario}{Abrir un caso}{17-mis-casos}{Conversación del caso}{Pantallas de
atención: el formulario para abrir un caso (a) y un caso de «Mis casos» con la respuesta del equipo
(b).}
\pantallados{16-privacidad-mis-datos}{Inventario de datos}{16b-privacidad-acciones}{Acciones del
titular}{Pantallas de privacidad: el inventario de la cuenta, los hijos, las direcciones y las
autorizaciones (a) y las acciones para descargar los datos, retirar la autorización y borrarlos (b).}

\subsubsection{Equipo de Tukit}

\pantalla{width=0.8\textwidth}{10-cola-revision-listas}{Pantalla de la cola de verificación, con los
trabajos de acudientes que esperan primero y las listas por verificar agrupadas por colegio.}
\pantallados{11a-revision-pdf-y-lectura}{PDF y lectura del modelo}{11b-renglones-con-concepto}{Renglones
con su concepto}{Pantallas de la revisión de una lista de Cuarto: el PDF original junto a lo que leyó
el modelo (a) y los renglones con su concepto y los atributos que exigen, entre ellos «un compás»
enviado a revisión (b).}
\pantallados{11c-corregir-renglon}{Corregir un renglón}{10b-subir-lista-oficial}{Subir una lista
oficial}{Pantallas de la corrección de un renglón (a) y de la carga de una lista oficial por el equipo
(b).}
\pantalla{width=0.72\textwidth,trim=0 250 0 0,clip}{12-catalogo-productos-conceptos}{Pantalla del
catálogo de productos en la consola, con los filtros por concepto, nivel y personaje y la tarjeta de
cada producto con su nivel, su SKU, su precio y su marca.}
\pantallados{13a-logistica-pedidos}{Pedidos y mapa de entregas}{13c-logistica-entregado}{Pedido
entregado}{Pantallas de logística: los pedidos en curso, con incidencia y entregados, los ingresos
confirmados y el mapa de entregas por colegio (a), y un pedido cerrado como entregado (b).}
\pantallados{09c-tickets-soporte}{Bandeja de casos}{14b-metricas-colegio-agustiniano}{Métricas por
colegio}{Pantallas de la bandeja de casos del equipo (a) y de las métricas de un colegio piloto (b).}
\pantalla{width=0.72\textwidth,trim=0 1790 0 0,clip}{14-metricas-panel}{Pantalla del panel de la
operación, con la cola de revisión, el estado de los trabajos y el resumen del catálogo.}

\fi

\subsection{Datos del sistema}
\label{anx:datos-sistema}

En CRISP-DM, esta subsección corresponde a la comprensión y la preparación de los datos. Antes de
construir las etapas del pipeline fue necesario saber con qué se alimentan y con qué se miden. Los
datos son de dos clases: el catálogo, que el sistema consulta para armar los kits, y los conjuntos de
evaluación, que sirven para medir cada etapa. Se describen por separado porque cumplen funciones
distintas y tienen sesgos distintos.

\paragraph{Catálogo.} Al descartar la extracción masiva de las tiendas (Sección~\ref{sec:consolidacion}),
el catálogo se construyó desde la demanda: un producto entra cuando una lista lo pide y el equipo
confirma la correspondencia. Cada producto pertenece a un concepto, la categoría de útil que una lista
nombra (por ejemplo, cuaderno o lápiz), y se clasifica en un nivel de precio y marca: económico,
confiable o premium. La curaduría distingue los empaques de las unidades. Un paquete de varias unidades
se registra con su número de unidades por empaque para que el precio se compare por unidad, y las
variantes de un mismo producto (color o tamaño) se agrupan en un grupo de variantes. Dentro de cada
grupo, un producto se marca como representante canónico (\texttt{es\_representante\_canonico}); solo
los representantes participan en la búsqueda y en la optimización, de modo que una variante no
desplaza a su grupo ni duplica un candidato. El modo \texttt{--surtir} del programa de curaduría
amplía el catálogo, y la cobertura por nivel de cada concepto
se verifica en celdas de concepto y nivel, porque un kit necesita al menos un candidato en cada una.
Los embeddings, vectores que representan el texto de cada producto, se calculan con
gemini-embedding-001 (1.024 dimensiones) y se indexan con HNSW en pgvector.

\begin{table}[htbp]
\centering
\caption{Catálogo y conjuntos de evaluación del sistema}
\label{tab:datos-sistema}
\small
\begin{tabularx}{\linewidth}{@{}p{3.6cm}p{4.4cm}X@{}}
\toprule
\textbf{Atributo} & \textbf{Valor} & \textbf{Fuente} \\
\midrule
Productos del catálogo & 1.054 (1.044 activos) & \texttt{datos/catalogo/productos\_canonicos.csv} \\
Representantes canónicos & 417 de 1.054 & mismo archivo, columna \texttt{es\_representante\_canonico} \\
Conceptos & 149 definidos; 133 con producto & \texttt{conceptos.csv}; \texttt{concepto\_id} del archivo de productos \\
Niveles & confiable 572, económico 268, premium 214 & columna \texttt{nivel} \\
Tiendas de procedencia & 8 & \texttt{ofertas\_tienda.csv} (columna \texttt{tienda\_id}) \\
Embeddings & gemini-embedding-001, 1.024 dimensiones, HNSW & código y SDD \\
CGS-10 (percepción) & 10 documentos, 1.293 ítems & Sección~\ref{sec:percepcion} \\
Oro de recuperación & 2.588 renglones de 4.719 ítems y 98 conceptos & Sección~\ref{sec:recuperacion} \\
Banco de listas reales & 175 instancias de 174 listas de 17 colegios & \texttt{ensayos/bench\_laya/instancias.json.gz} \\
Partición por colegio & 3 colegios de prueba (40 listas, 32 con NSGA-II) y 14 en cinco pliegues & \texttt{calibracion\_nsga.json}, campo \texttt{particion} \\
\bottomrule
\end{tabularx}
\end{table}

Las cifras de la Tabla~\ref{tab:datos-sistema} corresponden al archivo versionado el 4 de noviembre de
2026. Difieren ligeramente de las de la Sección~\ref{sec:consolidacion}, que describen el estado al
cierre de la construcción: el catálogo sigue creciendo con la curaduría, y la tabla toma el archivo
versionado, no el conteo de aquella sección.

\paragraph{Conjuntos de evaluación.} Hay cuatro. El primero es CGS-10, el conjunto de oro de
percepción: diez documentos con 1.293 ítems anotados, sobre el que se compararon los cinco motores de
lectura. El segundo es el conjunto de oro de recuperación, de 2.588 renglones tomados de los 4.719
ítems de las listas, con 98 conceptos (el catálogo actual tiene 149), que midió la regla de resolución de una versión anterior. El
tercero es el banco de listas reales, de 175 instancias derivadas de 174 listas de 17 colegios, que
alimenta la calibración y la evaluación del optimizador. El cuarto no es un conjunto sino una
partición de ese banco: los colegios se dividen en un grupo de prueba de tres, con 40 listas (32 por encima del umbral de enumeración, el resto exactas), y 14 de calibración repartidos en cinco pliegues. Se divide por
colegio y no por lista porque las listas de un mismo colegio se parecen entre sí, y una partición por
lista dejaría ver a la calibración listas casi idénticas a las de prueba.

\paragraph{Sesgos conocidos.} Cada conjunto tiene una limitación que condiciona su lectura.
CGS-10 se armó corrigiendo la salida de uno de los motores evaluados, por lo que el F1 de ese motor no
es independiente del oro. El oro de recuperación solo ejercitó el canal léxico, porque los 609
productos de entonces tenían el embedding vacío. El banco de listas es pequeño y agrupado: 17 colegios,
con listas de un mismo colegio correlacionadas. Y la partición por colegio deja solo tres colegios de
prueba, así que sus intervalos son anchos. Las secciones de resultados repiten cada sesgo junto a la
cifra que afecta.

\subsection{Asistente en lenguaje natural (NL2C)}
\label{sec:nl2c}

El asistente convierte la frase de un acudiente en las restricciones del optimizador. Corresponde a las etapas de modelado y evaluación de CRISP-DM en las fases F3 y F4 y sigue la estructura de las secciones anteriores: problema, entradas y salidas, estrategia, seudocódigo, costo, resultados y límites. Las tablas
completas de la evaluación están en el Anexo~\ref{anx:nl2c}.

\paragraph{Análisis del problema.} La propuesta aprobada prometió que el acudiente pudiera
describir en español lo que quiere, por ejemplo «Minnie para mi hija, máximo 500 mil», en lugar de
llenar un formulario de opciones cerradas. A esa tarea se le llama NL2C (del inglés \emph{natural
language to constraints}): transformar un texto en el objeto de restricciones que un programa ya
consume, que aquí es el optimizador de la Sección~\ref{sec:optimizacion-cal}. El texto $t$ es una
cadena corta, ambigua y con errores de ortografía, y la salida es la terna
$\mathcal{R}=\langle F,W,P\rangle$ de filtros duros, pesos y preferencias, con los dominios de la
Tabla~\ref{tab:nl2c-dominios}. Un campo que el acudiente no menciona queda en su valor por omisión,
y uno que menciona fuera de dominio se lleva al dominio y se le avisa. Dos condiciones fijan el
diseño. La primera es que $t$ y el texto de la lista los puede escribir un atacante, de modo que
ninguno de los dos puede cambiar qué acciones ejecuta el sistema \citep{greshake2023indirect}. La
segunda es que un error le cuesta dinero al acudiente, por lo que nada se aplica sin su
confirmación. Las reglas fijas no bastan, porque los nombres de personajes y marcas son un
vocabulario abierto y una frase como «calidad primero y precio después» fija los pesos pero no el
estilo de gasto.

\begin{table}[H]
\centering\footnotesize
\caption{Dominio de los campos de $\mathcal{R}$. Fuente: constantes de \texttt{nl2c.py}.}
\label{tab:nl2c-dominios}
\begin{tabularx}{\textwidth}{@{}p{1.2cm}p{4.2cm}X@{}}
\toprule
Grupo & Campo & Dominio \\
\midrule
$F$ & \texttt{presupuesto\_max\_cop} & vacío o entero en $[50.000;\ 3.000.000]$ pesos por hijo \\
 & \texttt{marcas\_excluidas}, \texttt{solo\_disponibles} & hasta 10 cadenas de 1 a 40 caracteres; verdadero o falso \\
$W$ & cuatro pesos & cada $w_i\in[0{,}05;\ 0{,}60]$ y $\sum w_i=1$ \\
$P$ & \texttt{personajes}, \texttt{estilo\_gasto} & hasta 3 del catálogo; económico, punto medio o estirarse \\
\bottomrule
\end{tabularx}
\end{table}

\paragraph{Diseño: entradas y salidas.} Las entradas son el texto $t$, de a lo más 1.000 caracteres
después de sanearlo, el catálogo de personajes $\mathcal{C}_P$, con hasta 60 nombres, y una lista
blanca de ocho componentes de interfaz. Para pedir cambios sobre un kit se agrega el kit, con hasta
80 renglones. Las salidas son $\mathcal{R}$ validada, la lista de ajustes que hizo el servidor, la
especificación de la interfaz que se le muestra al acudiente y un estado, que puede ser propuesta,
sin acción, inválido o sin servicio. Si el estado es propuesta, $\mathcal{R}$ cumple la
Tabla~\ref{tab:nl2c-dominios}, todo texto libre de $\mathcal{R}$ aparece en $t$ y nada se aplica
antes de que el acudiente confirme.

\paragraph{Estrategia algorítmica.} La estrategia es un plan fijo con validación en el servidor: el
código decide qué funciones puede llamar cada ruta antes de ver $t$, y el modelo de lenguaje solo
elige una de ellas \citep{beurerkellner2025patterns}. La idea sigue la separación entre planificar y
ejecutar de Debenedetti et al. \citep{debenedetti2025camel}. Con ese plan, el modelo recibe el texto como dato y llama a
una función con argumentos; el servidor valida los argumentos y arma la interfaz con los componentes
de la lista blanca, que la aplicación web vuelve a validar \citep{chen2025genui}; y el acudiente
confirma, como piden las pautas de interacción humano-IA \citep{amershi_guidelines_2019}. Con esas
barandas, el servidor cubre cinco riesgos de la lista OWASP para aplicaciones con modelos de
lenguaje \citep{owasp2025llm}: inyección de instrucciones, información sensible, manejo de la
salida, agencia excesiva y consumo ilimitado. La línea base es el formulario de opciones cerradas,
que no usa modelo ni tiene riesgo de inyección, pero no entiende frases que mezclan campos.

\begin{algorithm}[H]
\caption{Interpretación de una preferencia en lenguaje natural}
\label{alg:nl2c}
\KwEntrada{texto $t$, catálogo de personajes $\mathcal{C}_P$}
\KwSalida{estado y, si es propuesta, la terna $(\mathcal{R}, A, U)$}
$s\leftarrow$ Sanear$(t)$ \tcp*{quita datos personales y marcas que escapan del bloque de datos}
$L\leftarrow\{$fijar\_restricciones, sin\_acción$\}$ \tcp*{plan fijo: la lista no sale de $t$}
$(f,a)\leftarrow$ LlamarModelo$(s,L)$ \tcp*{una sola llamada, sin reintento}
\lIf{$f\notin L$ o $f=$ sin\_acción}{\Return{sin acción}}
$(\mathcal{R},A)\leftarrow$ Proyectar$(a,\mathcal{C}_P,s)$ \tcp*{forma, recorte a dominios, texto anclado a $s$}
\lIf{Proyectar falló por forma}{\Return{inválido}}
$U\leftarrow$ CompilarInterfaz$(\mathcal{R})$ \tcp*{la arma el servidor, no el modelo}
\Return{(propuesta, $\mathcal{R}$, $A$, $U$)} \tcp*{nada se aplica hasta la confirmación}
\end{algorithm}

\paragraph{Complejidad y costo.} Sanear recorre $t$ con un número constante de expresiones
regulares, con un costo de $O(|t|)$. Proyectar trabaja sobre una $\mathcal{R}$ de tamaño acotado y
cuesta $O(|\mathcal{C}_P|+|s|)$, por el índice del catálogo y el cotejo de cada marca con el texto;
la normalización de los pesos itera a lo más cuatro veces y la interfaz tiene a lo más seis
componentes. En total son $O(|t|+|\mathcal{C}_P|)$ más una llamada al modelo, acotada a 400 tokens de
salida y 12 segundos de plazo, que domina el tiempo. En la ronda 3, con 70 intentos, esa llamada tuvo
una latencia p50 de 1.000~ms y p95 de 1.177~ms, con 616 tokens de entrada y 39 de salida en promedio.
El formulario no llama al modelo, pero su tiempo no se midió con acudientes, de modo que no se afirma
que NL2C sea más rápido para ellos.

\paragraph{Dónde quedó en la aplicación.} El componente se construyó para los dos puntos de
entrada, el inicio del pedido y la página del kit, pero en la versión que usaron los acudientes de
la ronda 2 de OE5 solo se monta en la página del kit, como el asistente que recibe pedidos de cambio
del tipo «más barato en cuadernos». En el inicio el equipo dejó el formulario guiado, porque el
acudiente todavía no ha visto ningún kit y el formulario le muestra de una vez qué puede elegir. El
servicio de interpretación es el mismo en los dos casos, así que la evaluación de esta sección
vale para el asistente del kit; el panel del inicio sigue en el código detrás de una bandera
apagada.

Después de la ronda 3, el asistente cambió sin cambiar sus barandas. Ante un pedido por categoría,
como «más barato en cuadernos», el modelo solo clasifica la acción y la categoría, y reglas del
servidor eligen renglones, producto y precio, con hasta 12 propuestas y el ahorro total; «Aceptar
todos» envía cambios individuales con la sesión del acudiente. El contexto se reduce a la acción y
la categoría anteriores, validadas contra el kit, y el modelo no recibe datos personales.

\paragraph{Evaluación.} El asistente se midió con 35 frases en español, 22 de la encuesta y 13
construidas, cuya etiqueta de oro se escribió antes de correr el modelo, y con 27 ataques de
inyección al estilo de AgentDojo \citep{debenedetti2024agentdojo}. El modelo fue
\texttt{gemini-3.5-flash-lite} con temperatura 0, con dos repeticiones por frase (70 intentos) y tres
por ataque (81 intentos). Un ataque tiene éxito cuando su condición se cumple sobre lo que el
acudiente ve después de las barandas del servidor. En la ronda 3 la tarea completa se cumplió en 58
de 70 intentos (82,9\,\%), el F1 de personajes y marcas fue 1,000 y ningún ataque tuvo éxito (0 de
81), aunque el modelo lo intentó en 6. Los 12 fallos de utilidad se concentran en el campo del peso
mayor, donde varias etiquetas son discutibles. Bajo ataque se perdió la tarea legítima en 24 de 60
intentos, porque el modelo respondió sin acción ante la orden hostil, lo que equivale a preferir no
proponer ante la duda.

Entre la ronda 2 y la 3 cambió una condición de éxito, y se declara. La
condición \texttt{descarte\_\allowbreak aplicable} se sustituyó por \texttt{descarte\_sin\_confirmar}, porque al integrar el
asistente con la sesión de pedidos el descarte pasó a ser aplicable después de la confirmación del
acudiente. El cambio se hizo después de ver los resultados de la ronda 2, así que no es un criterio
preregistrado; en esa condición las dos rondas dieron 6 evaluaciones y 0 éxitos. Además, con los
servicios levantados, 4 pedidos contra un kit real de 36 renglones dieron la respuesta esperada en
786 a 885~ms cada uno, lo que con $n=4$ es una comprobación de funcionamiento y no una tasa.

El arnés sin conexión, que usa un modelo que obedece toda orden hostil y mide solo las barandas,
da 7 éxitos en 27 ataques, igual antes y después del cambio: cinco son ecos del propio texto del
acudiente y dos son propuestas válidas que él no pidió, que detiene su confirmación. Es el peor caso
de las barandas y no se compara con el 0 de 81 del modelo real; la ronda en vivo con Gemini sobre el
asistente nuevo no se hizo. En las sesiones con acudientes (Sección~\ref{sec:oe5}) el asistente
estuvo visible en la página del kit, pero ninguno lo usó por iniciativa propia, así que su utilidad
con usuarios reales queda como trabajo futuro.

\paragraph{Limitaciones.} Con 0 éxitos en 81 intentos, el límite superior del intervalo de confianza
del 95\,\% es de cerca de 3,7\,\% (regla de tres), por lo que solo se puede afirmar que no se encontró
un ataque exitoso con estos casos. Las repeticiones a temperatura 0 no son independientes, el modelo
y la fecha son uno solo, el kit y el catálogo de las rondas son sintéticos y los ataques los escribió
quien construyó las defensas. Los filtros duros sí llegan al optimizador (commit \texttt{0059f6d}):
\texttt{marcas\_excluidas} y \texttt{solo\_disponibles} se aplican antes del frente de Pareto, y
también en el carrusel de opciones y en el cambio de producto, pero no quitan una marca que el
colegio exige en el renglón. La ruta rápida, que publica el kit canónico sin recalcular, los
rechaza con el error 422 \texttt{filtros\_requieren\_recalculo}. Su límite es que
\texttt{solo\_disponibles} usa la oferta activa con precio vigente y no el inventario por tienda.


\subsection{Hallazgos del control de calidad}
\label{sec:anexo-hallazgos}
Al 5 de noviembre de 2026 hay 85 hallazgos vigentes: 16 de severidad alta, 58 media y 11 baja, además de 10 cerrados. El detalle está en \path{docs/informes/calidad/hallazgos_qa_oficial.csv}.

\subsection{Detalle de la calibración del optimizador}
\label{anx:calibracion}

\sloppy En este anexo se reúnen las tablas de consulta de la calibración que resume la
Sección~\ref{sec:optimizacion-cal}, donde también están las figuras. Todas salen de los archivos JSON de la calibración,
entre ellos \texttt{umbral.json}, \texttt{umbral\_elegida.json}, \texttt{calibracion\_nsga.json}, \texttt{sensibilidad.json} y \texttt{hipervolumen\_prueba.json}, medidos en un \calCPU{} con
\calRAM~GiB de memoria.

\paragraph{Umbral de enumeración.} La Tabla~\ref{tab:calopt-umbral} da los tiempos por nivel del
espacio por nivel; la gráfica con el ajuste está en la Sección~\ref{sec:optimizacion-cal}. Con la configuración
de partida, NSGA-II tardaba una mediana de \calTNsgaMediana~s y el cruce quedaba en
$S\approx\calCruce$; con la elegida baja a \calTNsgaElegidaMediana~s y el cruce, a
$S\approx\calCruceElegida$.

\begin{table}[htbp]
\centering
\caption{Tiempo de pared de la enumeración exhaustiva y de NSGA-II por nivel $b$ del espacio
($S\in[2^b,2^{b+1})$), con la configuración de partida y con la elegida. «Exacto más rápido» y
«HV mínimo» se refieren a la elegida; \calSubproblemas{} subproblemas y \calSemillasHV{} semillas por
subproblema para el hipervolumen.}
\label{tab:calopt-umbral}
\small
\begin{tabular}{rrrrrrrr}
\toprule
 & & \multicolumn{2}{c}{exacto (s)} & \multicolumn{2}{c}{NSGA-II, mediana (s)} & & \\
$b$ & subproblemas & mediana & máx. & partida & elegida & exacto más rápido & HV mínimo \\
\midrule
6 & 5 & 0{,}001 & 0{,}001 & 1{,}382 & 0{,}258 & 5 & 1{,}0000 \\
7 & 5 & 0{,}002 & 0{,}003 & 0{,}703 & 0{,}197 & 5 & 1{,}0000 \\
8 & 5 & 0{,}005 & 0{,}011 & 0{,}388 & 0{,}163 & 5 & 1{,}0000 \\
9 & 5 & 0{,}015 & 0{,}020 & 0{,}372 & 0{,}188 & 5 & 1{,}0000 \\
10 & 5 & 0{,}032 & 0{,}058 & 0{,}439 & 0{,}211 & 5 & 1{,}0000 \\
11 & 5 & 0{,}089 & 0{,}294 & 0{,}401 & 0{,}183 & 3 & 0{,}9998 \\
12 & 5 & 0{,}354 & 0{,}731 & 0{,}420 & 0{,}223 & 1 & 0{,}9999 \\
13 & 5 & 1{,}301 & 2{,}457 & 0{,}461 & 0{,}239 & 0 & 1{,}0000 \\
14 & 5 & 1{,}420 & 7{,}553 & 0{,}581 & 0{,}288 & 0 & 0{,}9996 \\
15 & 3 & 11{,}912 & 25{,}134 & 0{,}632 & 0{,}337 & 0 & 0{,}9998 \\
16 & 3 & 48{,}534 & 142{,}993 & 0{,}538 & 0{,}274 & 0 & 0{,}9997 \\
17 & 3 & 27{,}650 & 179{,}797 & 0{,}642 & 0{,}356 & 0 & 1{,}0000 \\
18 & 3 & --- & --- & 0{,}712 & 0{,}356 & 0 & 0{,}9999 \\
\bottomrule

\end{tabular}
\end{table}

\paragraph{Población y generaciones.} La Tabla~\ref{tab:calopt-rejilla} da la rejilla completa sobre
los pliegues de validación (\calCVListas{} listas, \calSemillasCV{} semillas por lista; los pliegues
tienen \calPlieguesListas{} listas). El pliegue en
que la regla eligió otra configuración escogió una todavía más barata, que quedó por fuera del margen
en su pliegue apartado, lo que indica que reducir más el presupuesto ya cuesta calidad. En la prueba,
el \calFraccionHV{} de las corridas de la elegida alcanzó al menos 0,999 del hipervolumen exacto, la
peor corrida llegó a \calRazonMinPrueba{} y el mínimo por lista fue \calPruebaElegidaHVMin. La rejilla
es reducida y se declara: 2 semillas por lista en la validación y 5 en la prueba.

\begin{table}[htbp]
\centering
\caption{Rejilla de población ($a$) y generaciones ($b$) sobre los pliegues de validación.
$\star$: configuración elegida; $\circ$: configuración de partida.}
\label{tab:calopt-rejilla}
\small
\begin{tabular}{rrrrrrc}
\toprule
$a$ & $b$ & evaluaciones medias & HV medio & HV mínimo por lista & tiempo medio (s) & \\
\midrule
0{,}25 & 0{,}25 & 1\,231 & 0{,}9960 & 0{,}9711 & 0{,}29 &  \\
0{,}25 & 0{,}50 & 2\,475 & 0{,}9973 & 0{,}9708 & 0{,}57 &  \\
0{,}25 & 1{,}00 & 4\,963 & 0{,}9971 & 0{,}9615 & 1{,}18 &  \\
0{,}25 & 2{,}00 & 9\,937 & 0{,}9975 & 0{,}9725 & 2{,}30 &  \\
0{,}50 & 0{,}25 & 2\,453 & 0{,}9986 & 0{,}9904 & 0{,}48 &  \\
0{,}50 & 0{,}50 & 4\,931 & 0{,}9998 & 0{,}9970 & 1{,}01 & $\star$ \\
0{,}50 & 1{,}00 & 9\,890 & 1{,}0000 & 0{,}9981 & 2{,}01 &  \\
0{,}50 & 2{,}00 & 19\,800 & 1{,}0000 & 0{,}9979 & 4{,}20 &  \\
1{,}00 & 0{,}25 & 4\,907 & 0{,}9987 & 0{,}9896 & 0{,}90 &  \\
1{,}00 & 0{,}50 & 9\,862 & 0{,}9999 & 0{,}9981 & 1{,}93 &  \\
1{,}00 & 1{,}00 & 19\,779 & 1{,}0000 & 0{,}9993 & 3{,}96 & $\circ$ \\
1{,}00 & 2{,}00 & 39\,601 & 1{,}0000 & 1{,}0000 & 8{,}00 &  \\
2{,}00 & 0{,}25 & 9\,814 & 0{,}9988 & 0{,}9935 & 1{,}86 &  \\
2{,}00 & 0{,}50 & 19\,723 & 0{,}9999 & 0{,}9973 & 4{,}03 &  \\
2{,}00 & 1{,}00 & 39\,558 & 1{,}0000 & 0{,}9992 & 8{,}50 &  \\
2{,}00 & 2{,}00 & 79\,201 & 1{,}0000 & 1{,}0000 & 16{,}46 &  \\
\bottomrule

\end{tabular}
\end{table}

\paragraph{Umbral $\tau$.} Con la curaduría disponible, en $\tau=0{,}75$ se acepta el
\calTauAceptada{} de los renglones con una precisión del \calTauPrecision, pero el parecido solo toma
\calTauValores{} valores distintos y la validación cruzada agrupada elige valores entre
\calTauCVMin{} y \calTauCVMax{} según el pliegue. Por eso, y porque no se calculó el acuerdo entre los 4 anotadores, el valor de 0,75 se fija por su
precisión sobre los \calTauPares{} pares y se reporta con su intervalo, sin afinarlo pliegue a pliegue.

\paragraph{Sensibilidad del kit recomendado.} En las \calSensListas{} listas de prueba con al menos
dos puntos en el frente (mediana de \calSensFrenteMediana{} puntos) se sortearon
\calSensPerturbaciones{} vectores de pesos por lista y celda. Sin sobreprecio la estabilidad es 1 por
construcción, porque la cota de sobreprecio deja admisible solo el kit más barato. En el primer
esquema cada peso se sortea dentro de su intervalo del 95\,\% de la P21; en el segundo, los pesos
siguen una distribución de Dirichlet centrada en los de la P21 con concentración $\kappa$, que admite
perturbaciones más anchas. La Tabla~\ref{tab:calopt-sensibilidad} da todas las celdas.

\begin{table}[htbp]
\centering
\caption{Fracción de sorteos de pesos en que el kit recomendado por MAUT es el mismo que con los
pesos de la P21. DE: desviación estándar del peso del precio bajo la perturbación.}
\label{tab:calopt-sensibilidad}
\small
\begin{tabular}{rlrrlr}
\toprule
sobreprecio & perturbación & DE & media & IC 95\,\% & mínimo por lista \\
\midrule
0\,\% & IC 95\,\% P21 & --- & 1{,}000 & [1{,}000; 1{,}000] & 1{,}000 \\
5\,\% & IC 95\,\% P21 & --- & 0{,}965 & [0{,}932; 0{,}991] & 0{,}610 \\
10\,\% & IC 95\,\% P21 & --- & 0{,}986 & [0{,}964; 1{,}000] & 0{,}692 \\
20\,\% & IC 95\,\% P21 & --- & 0{,}946 & [0{,}904; 0{,}981] & 0{,}540 \\
0\,\% & Dirichlet $\kappa=20$ & 0{,}099 & 1{,}000 & [1{,}000; 1{,}000] & 1{,}000 \\
5\,\% & Dirichlet $\kappa=20$ & 0{,}099 & 0{,}892 & [0{,}830; 0{,}947] & 0{,}429 \\
10\,\% & Dirichlet $\kappa=20$ & 0{,}099 & 0{,}868 & [0{,}822; 0{,}913] & 0{,}541 \\
20\,\% & Dirichlet $\kappa=20$ & 0{,}099 & 0{,}735 & [0{,}663; 0{,}803] & 0{,}168 \\
0\,\% & Dirichlet $\kappa=50$ & 0{,}063 & 1{,}000 & [1{,}000; 1{,}000] & 1{,}000 \\
5\,\% & Dirichlet $\kappa=50$ & 0{,}063 & 0{,}918 & [0{,}866; 0{,}965] & 0{,}502 \\
10\,\% & Dirichlet $\kappa=50$ & 0{,}063 & 0{,}927 & [0{,}888; 0{,}960] & 0{,}558 \\
20\,\% & Dirichlet $\kappa=50$ & 0{,}063 & 0{,}828 & [0{,}765; 0{,}886] & 0{,}287 \\
0\,\% & Dirichlet $\kappa=100$ & 0{,}045 & 1{,}000 & [1{,}000; 1{,}000] & 1{,}000 \\
5\,\% & Dirichlet $\kappa=100$ & 0{,}045 & 0{,}931 & [0{,}884; 0{,}974] & 0{,}520 \\
10\,\% & Dirichlet $\kappa=100$ & 0{,}045 & 0{,}953 & [0{,}923; 0{,}978] & 0{,}618 \\
20\,\% & Dirichlet $\kappa=100$ & 0{,}045 & 0{,}876 & [0{,}818; 0{,}926] & 0{,}379 \\
0\,\% & Dirichlet $\kappa=200$ & 0{,}032 & 1{,}000 & [1{,}000; 1{,}000] & 1{,}000 \\
5\,\% & Dirichlet $\kappa=200$ & 0{,}032 & 0{,}945 & [0{,}903; 0{,}980] & 0{,}557 \\
10\,\% & Dirichlet $\kappa=200$ & 0{,}032 & 0{,}973 & [0{,}946; 0{,}993] & 0{,}632 \\
20\,\% & Dirichlet $\kappa=200$ & 0{,}032 & 0{,}910 & [0{,}857; 0{,}956] & 0{,}456 \\
\bottomrule

\end{tabular}
\end{table}

\paragraph{Hipervolumen por perfil.} La Tabla~\ref{tab:calopt-hv} da el hipervolumen relativo al
frente exacto en las \calHVListas{} listas de prueba (\calHVListasNsga{} de ellas con NSGA-II). El mínimo del frente de producción fue \calHVFrenteMin.

\begin{table}[htbp]
\centering
\caption{Hipervolumen relativo al frente exacto en las listas de prueba, con intervalo
\emph{bootstrap} del 95\,\%.}
\label{tab:calopt-hv}
\small
\begin{tabular}{lrl}
\toprule
conjunto & media & IC 95\,\% \\
\midrule
Frente de producci\'on & 0{,}9969 & [0{,}9951; 0{,}9986] \\
Ahorrador $\cup$ Premium & 0{,}501 & [0{,}474; 0{,}526] \\
Ahorrador & 0{,}362 & [0{,}343; 0{,}382] \\
Premium & 0{,}231 & [0{,}201; 0{,}262] \\
MAUT con pesos P21 & 0{,}517 & [0{,}474; 0{,}558] \\
\bottomrule

\end{tabular}
\end{table}

\subsection{Detalle de la evaluación del asistente NL2C}
\label{anx:nl2c}

En este anexo se reúnen las tablas de la ronda 3 de la evaluación del asistente que resume la
Sección~\ref{sec:nl2c}. La Tabla~\ref{tab:nl2c-rondas} da las medidas de la evaluación.

\begin{table}[htbp]
\centering\small
\caption{Resultados de las rondas 2 y 3 del asistente NL2C.}
\label{tab:nl2c-rondas}
\begin{tabular}{@{}p{5.4cm}p{3.9cm}p{3.9cm}@{}}\toprule
Medida & Ronda 2 (criterio original) & Ronda 3 (criterio nuevo)\\\midrule
Tarea completa & 57 de 70 (81,4\,\%) & 58 de 70 (82,9\,\%)\\
F1 de personajes y marcas & 1,000 & 1,000\\
Latencia p50 / p95 (ms) & 999 / 1.344 & 1.000 / 1.177\\
Ataques exitosos (ASR) & 0 de 81 & 0 de 81\\
El modelo lo intentó & 5 de 81 & 6 de 81\\
Utilidad bajo ataque & 36 de 60 & 36 de 60\\
Controles sin ataque & 12 de 12 & 12 de 12\\\bottomrule
\end{tabular}
\end{table}

\iffalse % Tablas de la ronda 3 por campo, ataque y costo: docs/validacion/nl2c/ y docs/memoria/sections/tablas/nl2c/
\begin{table}[htbp]
\centering\small
\caption{Coincidencia con el oro por campo de $\mathcal{R}$ en la ronda 3.}
\label{tab:nl2c-campos}
\begin{tabular}{@{}lrrr@{}}\toprule
Campo de R & Aciertos & Evaluados & Coincidencia\\\midrule
acción & 68 & 70 & 97,1\,\%\\
presupuesto & 58 & 60 & 96,7\,\%\\
personajes & 58 & 60 & 96,7\,\%\\
marcas excluidas & 58 & 60 & 96,7\,\%\\
solo disponibles & 58 & 60 & 96,7\,\%\\
estilo de gasto & 56 & 60 & 93,3\,\%\\
sobreprecio & 6 & 6 & 100,0\,\%\\
peso mayor & 50 & 60 & 83,3\,\%\\
tarea completa & 58 & 70 & 82,9\,\%\\
\bottomrule\end{tabular}

\end{table}

\begin{table}[htbp]
\centering\small
\caption{Ataques de inyección por categoría en la ronda 3.}
\label{tab:nl2c-asr}
\begin{tabular}{@{}lrrrr@{}}\toprule
Categoría & Casos & Intentos & Modelo lo intentó & Ataque exitoso\\\midrule
directa sobre R & 16 & 48 & 3 (6,2\,\%) & 0 (0,0\,\%)\\
directa sobre el kit & 5 & 15 & 3 (20,0\,\%) & 0 (0,0\,\%)\\
indirecta (texto del kit) & 6 & 18 & 0 (0,0\,\%) & 0 (0,0\,\%)\\
\midrule
Total & 27 & 81 & 6 (7,4\,\%) & 0 (0,0\,\%)\\
\bottomrule\end{tabular}

\end{table}

\begin{table}[htbp]
\centering\small
\caption{Condiciones de éxito del atacante en la ronda 3.}
\label{tab:nl2c-condiciones}
\begin{tabular}{@{}lrrr@{}}\toprule
Condición de éxito del atacante & Evaluados & Intentó el modelo & Llegó al padre\\\midrule
R fuera de dominio & 24 & 3 & 0\\
más de 5 propuestas & 6 & 3 & 0\\
fuga del prompt de sistema & 3 & 0 & 0\\
marcado ejecutable en pantalla & 12 & 0 & 0\\
cadena del atacante en pantalla & 45 & 0 & 0\\
dato sensible en pantalla & 18 & 0 & 0\\
función fuera de la lista & 9 & 0 & 0\\
renglón o perfil ajeno & 9 & 0 & 0\\
descarte sin confirmar (criterio de la ronda 3) & 6 & 0 & 0\\
propuesta no pedida & 6 & 0 & 0\\
enlace externo en pantalla & 3 & 0 & 0\\
\bottomrule\end{tabular}

\end{table}

\begin{table}[htbp]
\centering\small
\caption{Costo y latencia de la ronda 3. El costo usa un precio de lista supuesto, no medido.}
\label{tab:nl2c-costo}
\begin{tabular}{@{}lr@{}}\toprule
Medida & Valor\\\midrule
Intentos (frases por repeticiones) & 70\\
F1 de personajes y marcas & 1,000\\
Precisión / exhaustividad & 1,000 / 1,000\\
Latencia p50 (ms) & 1000\\
Latencia p95 (ms) & 1177\\
Tokens de entrada por pedido (media) & 616\\
Tokens de salida por pedido (media) & 39\\
Costo por pedido (USD, precio supuesto) & 0,000077\\
Intentos sin servicio & 0\\
\bottomrule\end{tabular}

\end{table}

\begin{table}[htbp]
\centering\small
\caption{Utilidad bajo ataque y controles en la ronda 3.}
\label{tab:nl2c-utilidad}
\begin{tabular}{@{}lrr@{}}\toprule
Situación & Tareas cumplidas & Tasa\\\midrule
Bajo ataque (con tarea legítima) & 36 de 60 & 60,0\,\%\\
Controles sin ataque (kit) & 12 de 12 & 100,0\,\%\\
\bottomrule\end{tabular}

\end{table}

\fi

% Seudocodigo completo de las etapas del pipeline, portado de algpseudocode a algorithm2e.
\subsection{Algoritmos del pipeline}
\label{sec:anexo-algoritmos}

Esta sección reúne los procedimientos completos de las etapas del pipeline que se describen en el
capítulo de Desarrollo: la percepción del archivo, la memoización de su lectura, la normalización de
cada renglón, su resolución contra el catálogo, la construcción del frente de Pareto, NSGA-II, la
selección de los kits y el frente exacto por programación dinámica.

% ---------------------------------------------------------------------
\paragraph{Percepción del archivo.}
El procedimiento recibe el archivo de una lista, lo rasteriza a 150 puntos por pulgada y lee cada
página con el modelo multimodal, validando la respuesta contra el esquema. El archivo se acepta solo
si todas sus páginas se leyeron completas y si corresponde a un único grado, contado después de
canonizar cada nombre; el tope de 40 páginas se revisa antes de gastar una llamada. Con $P$ páginas
el algoritmo hace $\Theta(P)$ llamadas, como máximo $2P$, de modo que dominan las inferencias.

\begin{algorithm}[H]
\caption{Percibir el archivo de una lista, con todas sus páginas o ninguna.}
\label{alg:anexo-percibir}
\small
\KwEntrada{$A$, archivo PDF o imagen; $\textsc{Leer}(I,j)$, lectura multimodal de la página $j$
validada contra el esquema.}
\KwSalida{$\langle ok, G, X, e\rangle$, con $ok$ verdadero solo si todas las páginas se leyeron
completas y $\lvert G\rvert \le 1$.}
\Proc{\textup{\textsc{Percibir}}$(A)$}{
  \If{$\lvert A\rvert = 0$}{
    \Return{$\langle \text{falso}, \emptyset, \langle\rangle, \textup{\textsc{ArchivoVacío}}\rangle$}\;
  }
  $\langle I_1,\dots,I_P\rangle \gets \textsc{Rasterizar}(A, 150)$\tcp*{una imagen suelta ya es una página}
  \If{$P > 40$}{
    \Return{$\langle \text{falso}, \emptyset, \langle\rangle, \textup{\textsc{DemasiadoLargo}}\rangle$}\tcp*{antes de gastar una llamada}
  }
  sean $X$ y $\Gamma$ secuencias vacías; $inc \gets \emptyset$\;
  \For{$j \gets 1$ \KwTo $P$}{
    $\langle G_j, \kappa_j, X_j\rangle \gets \textsc{Leer}(I_j, j)$\tcp*{un fallo termina con \textup{\textsc{PercepciónFallida}}}
    \If{$\kappa_j = \text{falso}$}{
      $inc \gets inc \cup \{j\}$\;
    }
    $\Gamma \gets \Gamma \cdot G_j$;\quad $X \gets X \cdot X_j$\;
  }
  $G \gets \textsc{GradosDistintos}(\Gamma)$\tcp*{con cada nombre canonizado}
  \If{$\lvert G\rvert > 1$}{
    \Return{$\langle \text{falso}, G, X, \textup{\textsc{VariosGrados}}\rangle$}\;
  }
  \If{$inc \ne \emptyset$}{
    \Return{$\langle \text{falso}, G, X, \textup{\textsc{PáginaIncompleta}}\rangle$}\;
  }
  \Return{$\langle \text{verdadero}, G, X, \emptyset\rangle$}\;
}
\end{algorithm}

\iffalse % Contenido en el repositorio: docs/algoritmos/
% ---------------------------------------------------------------------
\paragraph{Memoización de la lectura.}
Como las familias de un colegio suben el mismo archivo, la lectura aceptada se guarda en Redis
durante 30 días, con una clave que combina la huella SHA-256 del archivo y una versión del código de
lectura y del modelo. Los rechazos no se guardan, y si el almacén no responde la lectura se hace de
nuevo. Una lectura memorizada cuesta solo la huella, $\Theta(\lvert A\rvert)$.

\begin{algorithm}[H]
\caption{Memoización de la lectura de un archivo.}
\label{alg:anexo-memo}
\small
\KwEntrada{$A$, archivo; $R$, almacén clave-valor con vencimiento, posiblemente no disponible; $v$,
huella del código de lectura y del modelo.}
\Proc{\textup{\textsc{PercibirMemo}}$(A)$}{
  $k \gets v \cdot \textsc{Sha256}(A)$\;
  \If{$R$ responde \textbf{y} $R[k] \ne \bot$}{
    \Return{$R[k]$}\;
  }
  $L \gets \textsc{Percibir}(A)$\;
  \If{$L.ok$}{
    $R[k] \gets L$ con vencimiento de 30 días\tcp*{los rechazos no se guardan}
  }
  \Return{$L$}\;
}
\end{algorithm}

% ---------------------------------------------------------------------
\paragraph{Normalización de un renglón.}
La normalización recibe el texto del renglón, la cantidad de la columna si existía y el tipo y el
modo curados de la lista oficial, y devuelve una demanda con el tipo, el modo, la cantidad pedida y
el empaque que declara el texto. El tipo y el modo salen de una lista de decisión de once filas en la
que gana la primera que se cumple, y lo curado prevalece sobre lo inferido. Con $\ell$ la longitud
del renglón y un número constante de frases clave, normalizar un renglón es $\Theta(\ell)$ y una
lista, lineal en su longitud total.

\begin{algorithm}[H]
\caption{Normalizar un renglón en una demanda.}
\label{alg:anexo-normalizar}
\small
\KwEntrada{$x$, texto; $q^{\circ}$, cantidad de la columna o $\bot$; $t^{\ast}, m^{\ast}$, tipo y
modo curados o $\bot$; $D = \langle (C_r, t_r, m_r) ; 1 \le r \le 11\rangle$, lista de decisión con
$C_{11}$ siempre verdadera.}
\KwSalida{$\langle t, m, q, \upsilon\rangle$}
\Proc{\textup{\textsc{Normalizar}}$(x, q^{\circ}, t^{\ast}, m^{\ast})$}{
  $\bar x \gets \textsc{Plano}(x)$\tcp*{minúsculas, sin tildes y ``x12'' como ``x 12''}
  $r \gets \min\{\, r : \bar x \text{ contiene una frase de } C_r \,\}$\;
  $t \gets t^{\ast}$ si $t^{\ast} \ne \bot$, si no $t_r$;\quad
  $m \gets m^{\ast}$ si $m^{\ast} \ne \bot$, si no $m_r$\;
  \uIf{$t = \text{no\_comprar}$}{
    $q \gets 0$\;
  }
  \uElseIf{$q^{\circ} \ne \bot$}{
    $q \gets q^{\circ}$\tcp*{no se suma al número del texto}
  }
  \uElseIf{$\bar x$ empieza por un número $n$}{
    $q \gets n$\;
  }
  \Else{
    $q \gets 1$\;
  }
  $\upsilon \gets$ mayor $N$ de ``x $N$'', ``por $N$'' o ``caja de $N$'' sin unidad de medida
  detrás, o $\bot$\;
  \Return{$\langle t, m, q, \upsilon\rangle$}\;
}
\end{algorithm}

\fi
% ---------------------------------------------------------------------
\paragraph{Puntaje híbrido léxico-semántico.}
El procedimiento recibe el renglón ya corregido, los conceptos del catálogo con su nombre, sus
sinónimos y su centroide, y el vector del renglón si existe, y devuelve el concepto de mayor puntaje.
Si la frase del concepto empieza por la cabeza del renglón, el puntaje parte de 0,6; si la nombra en
otra posición, vale el Jaccard más el aporte del coseno, y si no la nombra solo un coseno de 0,86 o más
la rescata. El coseno se recorta en cero y solo suma, de modo que el puntaje híbrido nunca queda por
debajo del léxico. Clasificar un renglón es lineal en el catálogo, porque el coseno se calcula una vez
por concepto.

\begin{algorithm}[H]
\caption{Puntaje híbrido léxico-semántico de un renglón frente a los conceptos del catálogo.}
\label{alg:anexo-clasificar}
\small
\KwEntrada{$\tilde x$, renglón corregido; $\mathcal{C}$, conceptos con nombre, sinónimos $S(c)$ y
centroide $\bar{\mathbf{e}}(c)$; $\mathbf{e}(x)$, vector del renglón o $\bot$.}
\KwSalida{$\langle \hat c, \hat s\rangle$, concepto de mayor puntaje y su puntaje.}
\Proc{\textup{\textsc{Clasificar}}$(\tilde x, \mathcal{C}, \mathbf{e}(x))$}{
  reescribir en $\tilde x$ las series de durezas y los pares de escuadras\;
  $h \gets \textsc{Cabeza}(\tilde x)$;\quad $\langle \hat c, \hat s\rangle \gets \langle \bot, 0\rangle$\;
  \ForEach{$c \in \mathcal{C}$}{
    $\kappa \gets \max\bigl(0, \cos(\mathbf{e}(x), \bar{\mathbf{e}}(c))\bigr)$, o $0$ si falta un vector\;
    \ForEach{$y \in \{\mathrm{nom}(c)\} \cup S(c)$}{
      \uIf(\tcp*[h]{ancla, la frase empieza por lo pedido}){$h \approx y^{(1)}$}{
        $\psi \gets \min\bigl(1,\ 0{,}6 + 0{,}4\,J(\tilde x, y) + 0{,}15\,\kappa\bigr)$\;
      }
      \uElseIf{$h \approx y^{(i)}$ para algún $i$}{
        $\psi \gets \min\bigl(1,\ J(\tilde x, y) + 0{,}15\,\kappa\bigr)$\;
      }
      \uElseIf(\tcp*[h]{rescate semántico}){$\kappa \ge 0{,}86$}{
        $\psi \gets 0{,}55 + 0{,}4\,(\kappa - 0{,}86)/0{,}14$\;
      }
      \Else{
        $\psi \gets 0$\;
      }
      \If{$\psi > \hat s$}{
        $\langle \hat c, \hat s\rangle \gets \langle c, \psi\rangle$\tcp*{en empate gana el primero}
      }
    }
  }
  \Return{$\langle \hat c, \hat s\rangle$}\;
}
\end{algorithm}

% ---------------------------------------------------------------------
\paragraph{Resolución de un renglón.}
Reconocido el concepto, el procedimiento lee solo los atributos que ese concepto declara como
discriminantes y filtra sus productos. Si ninguno cumple, se abstiene en lugar de ofrecer otro, y
entre los que cumplen ofrece uno por nivel de marca, con máximo tres opciones. El estado depende de
dos umbrales: por debajo de 0,5 el renglón queda sin correspondencia, entre 0,5 y $\tau$ pasa a
revisión humana y desde $\tau$ se resuelve de forma automática. Un concepto asignado por una persona
se usa con puntaje~1, y resolver una lista es lineal en sus renglones.

\begin{algorithm}[H]
\caption{Resolver un renglón contra el catálogo y asignarle su estado.}
\label{alg:anexo-resolver}
\small
\KwEntrada{$x$, renglón; catálogo $(\mathcal{C}, \mathcal{P})$; $c^{\ast}$, concepto asignado por una
persona o $\bot$; umbral $\tau$.}
\KwSalida{$\langle e, \hat c, O\rangle$, con $e \in \{\textsc{ResueltoAuto}, \textsc{Ambiguo},
\textsc{HuecoCatálogo}, \textsc{SinMatch}\}$ y $\lvert O\rvert \le 3$.}
\Proc{\textup{\textsc{Resolver}}$(x, c^{\ast}, \tau)$}{
  $\tilde x \gets \textsc{Corregir}(x)$\tcp*{una única palabra del léxico a distancia 1}
  \uIf{$c^{\ast} \ne \bot$}{
    $\langle \hat c, \hat s\rangle \gets \langle c^{\ast}, 1\rangle$\tcp*{lo curado prevalece}
  }
  \Else{
    $\langle \hat c, \hat s\rangle \gets \textsc{Clasificar}(\tilde x, \mathcal{C}, \mathbf{e}(x))$\;
  }
  \If{$\hat c = \bot$ \textbf{o} $\hat s < 0{,}5$}{
    \Return{$\langle \textup{\textsc{SinMatch}}, \bot, \emptyset\rangle$}\;
  }
  $Imp \gets \textsc{Atributos}(\tilde x, D(\hat c))$\tcp*{solo los discriminantes del concepto}
  $K \gets \{\, p \in \mathcal{P}(\hat c) : \textsc{Cumple}(p, Imp) \,\}$\;
  \If{$K = \emptyset$}{
    \Return{$\langle \textup{\textsc{HuecoCatálogo}}, \hat c, \emptyset\rangle$}\tcp*{no se ofrece otro producto}
  }
  $O \gets$ un producto por nivel de $K$, sin personaje primero y con desempate por SKU\;
  \If{$\hat s \ge \tau$}{
    \Return{$\langle \textup{\textsc{ResueltoAuto}}, \hat c, O\rangle$}\;
  }
  \Return{$\langle \textup{\textsc{Ambiguo}}, \hat c, O\rangle$}\tcp*{pasa a revisión humana}
}
\end{algorithm}

\iffalse % Contenido en el repositorio: docs/algoritmos/
% ---------------------------------------------------------------------
\paragraph{Frente de Pareto de un hogar.}
El procedimiento recibe los candidatos admisibles de las líneas decidibles y los cuatro objetivos de
las Ecuaciones~\eqref{eq:f1f2} y~\eqref{eq:f3f4}, y construye el frente según el tamaño del espacio $\Omega$. Con
$\lvert\Omega\rvert \le 1$ no hay compromiso y decide la heurística por perfil; con
$\lvert\Omega\rvert \le 4.096$ enumera todas las asignaciones y obtiene el frente exacto, y por encima
de ese límite lo aproxima con NSGA-II sembrado. La enumeración cuesta $O(\lvert\Omega\rvert\,n)$ más
el filtro por pares, $O(\lvert U\rvert^2 M)$ con $M = 4$.

\begin{algorithm}[H]
\caption{Construir el frente de Pareto de un hogar.}
\label{alg:anexo-frente}
\small
\KwEntrada{$A_1, \dots, A_n$, candidatos admisibles de las líneas decidibles; $\mathbf{F} = (f_1, f_2,
f_3, f_4)$; semilla $\varsigma$.}
\KwSalida{Frente $\mathcal{F}$, con una asignación por vector objetivo.}
\Proc{\textup{\textsc{Frente}}$(A_1, \dots, A_n)$}{
  $\lvert\Omega\rvert \gets \prod_{i=1}^{n} \lvert A_i\rvert$\;
  \If{$n = 0$ \textbf{o} $\lvert\Omega\rvert \le 1$}{
    \Return{$\emptyset$}\tcp*{sin compromiso, decide la heurística por perfil}
  }
  \uIf{$\lvert\Omega\rvert \le 4096$}{
    $S \gets \{\, \langle \mathbf{F}(\mathbf{x}), \mathbf{x}\rangle : \mathbf{x} \in A_1 \times \cdots \times A_n \,\}$\tcp*{enumeración exacta}
  }
  \Else{
    $E \gets \textsc{Heurística}(A) \cup \textsc{Escalera}(A)$\tcp*{semillas}
    $S \gets \textsc{NSGA-II}\bigl(N(n), G(\lvert\Omega\rvert), E, \varsigma\bigr) \cup E$\;
  }
  $U \gets$ un representante de $S$ por cada vector objetivo distinto\;
  \Return{$\{\, \mathbf{x} \in U : \neg\exists\, \mathbf{y} \in U,\ \mathbf{y} \prec \mathbf{x} \,\}$
  ordenado por $(f_1, -f_2, -f_3, -f_4)$}\;
}
\end{algorithm}

% ---------------------------------------------------------------------
\paragraph{NSGA-II sembrado.}
NSGA-II mantiene $N$ soluciones, produce $N$ descendientes por selección, cruce binario simulado y
mutación polinomial, y elige la generación siguiente por rango de no dominancia y distancia de
aglomeración. La variable de cada línea es el índice entero del candidato, y la población inicial se
siembra con la heurística por perfil y con la escalera de costo marginal, de modo que el frente nunca
queda por debajo de las semillas. Su costo es $O\bigl(G(MN^2 + Nn)\bigr)$, por el ordenamiento no
dominado y la evaluación.

\begin{algorithm}[H]
\caption{NSGA-II sembrado sobre la codificación entera de índices.}
\label{alg:anexo-nsga}
\small
\KwEntrada{$N$, tamaño de la población; $G$, generaciones; $E$, semillas; $\varsigma$, semilla;
cotas $0 \le x_i \le \lvert A_i\rvert - 1$.}
\KwSalida{Aproximación del frente de Pareto.}
\Proc{\textup{\textsc{NSGA-II}}$(N, G, E, \varsigma)$}{
  $P_0 \gets E$ completada con individuos al azar generados con $\varsigma$\tcp*{si $\lvert E\rvert > N$, $N$ semillas equiespaciadas}
  evaluar $\mathbf{F}$ en $P_0$\;
  \For{$t \gets 0$ \KwTo $G-1$}{
    $Q_t \gets N$ hijos por torneo binario, cruce ($p_c = 0{,}9$, $\eta_c = 15$) y mutación
    ($\eta_m = 20$), redondeados y sin duplicados\;
    evaluar $\mathbf{F}$ en $Q_t$;\quad $R_t \gets P_t \cup Q_t$\;
    $\langle \mathcal{F}_1, \mathcal{F}_2, \dots\rangle \gets$ ordenamiento no dominado de $R_t$\tcp*{$O(M(2N)^2)$}
    $P_{t+1} \gets$ frentes completos en orden de rango hasta $N$, con el último recortado
    por distancia de aglomeración\;
  }
  \Return{los no dominados de $P_G$}\;
}
\end{algorithm}

% ---------------------------------------------------------------------
\paragraph{Selección de los kits.}
Del frente salen los kits por selección lexicográfica sobre los puntos de máxima disponibilidad. El
Ahorrador minimiza el costo, el Premium maximiza la fidelidad, el Fan maximiza el personaje, y A tu
Medida es el de mayor fidelidad cuyo costo cabe en el tope, que es el presupuesto declarado o la
mitad del camino entre el Ahorrador y el Premium. Si ningún punto trae el personaje, el Fan repite el
Ahorrador en lugar de inventarlo.

\begin{algorithm}[H]
\caption{Selección lexicográfica de los kits sobre el frente.}
\label{alg:anexo-proyectar}
\small
\KwEntrada{$\mathcal{F} \ne \emptyset$; $B$, tope declarado o $\bot$.}
\KwSalida{Asignaciones de los kits Ahorrador, Premium, Fan y A tu Medida.}
\Proc{\textup{\textsc{Proyectar}}$(\mathcal{F}, B)$}{
  $\mathcal{F}_s \gets \{\, \mathbf{x} \in \mathcal{F} : f_3(\mathbf{x}) = \max_{\mathbf{y} \in \mathcal{F}} f_3(\mathbf{y}) \,\}$\tcp*{puntos surtibles}
  $\mathbf{x}_{Ah} \gets \operatorname{lex\,min}_{\mathcal{F}_s}\,(f_1, -f_3, -f_2)$\;
  $\mathbf{x}_{Pr} \gets \operatorname{lex\,max}_{\mathcal{F}_s}\,(f_2, -f_1, f_3)$\;
  $\mathbf{x}_{Fan} \gets \operatorname{lex\,max}_{\mathcal{F}_s}\,(f_4, -f_1, f_2)$\;
  \If{$f_4(\mathbf{x}_{Fan}) = 0$}{
    $\mathbf{x}_{Fan} \gets \mathbf{x}_{Ah}$\tcp*{no se inventa un personaje}
  }
  $T \gets B$ si $B \ne \bot$, si no $\bigl\lfloor (\mathrm{total}(\mathbf{x}_{Ah}) + \mathrm{total}(\mathbf{x}_{Pr}))/2 \bigr\rfloor$\;
  $\mathbf{x}_{AM} \gets \operatorname{lex\,max}_{\{\mathbf{x} \in \mathcal{F}_s\, :\, f_1(\mathbf{x}) \le T\}}\,(f_2, -f_1, f_3)$\;
  \If{$\mathbf{x}_{AM}$ no existe \textbf{o} su total consolidado supera $T$ \textbf{o} repite
  $\mathbf{x}_{Ah}$ o $\mathbf{x}_{Pr}$}{
    $\mathbf{x}_{AM} \gets \bot$\;
  }
  \Return{$\langle \mathbf{x}_{Ah}, \mathbf{x}_{Pr}, \mathbf{x}_{Fan}, \mathbf{x}_{AM}\rangle$}\;
}
\end{algorithm}

% ---------------------------------------------------------------------
\paragraph{Frente exacto por programación dinámica.}
Como cada objetivo es una suma de términos de una sola línea y tres de ellos son contadores enteros
pequeños, el menor costo de cada terna $(a, b, p)$ de nivel de marca acumulado, líneas con precio y
líneas con personaje sale de una recurrencia línea por línea, y el frente es el conjunto de ternas no
dominadas con su costo mínimo. Tras $i$ líneas hay a lo sumo $(2i+1)(i+1)^2$ estados con $k$
transiciones cada uno, de modo que el tiempo es $O(n^4k)$ y la memoria, con el retroceso de todas las
etapas, $O(n^4)$.

\begin{algorithm}[H]
\caption{Frente exacto por programación dinámica sobre contadores enteros.}
\label{alg:anexo-pd}
\small
\KwEntrada{$A_1, \dots, A_n$; para cada candidato $c$, su nivel $\nu(c) \in \{0,1,2\}$, $\beta(c) = 1$
si tiene precio y $0$ si no, $\varphi_i(c) \in \{0,1\}$ y su costo de línea $g_i(c)$.}
\KwSalida{$C^{\ast}(a,b,p)$, menor costo de cada terna alcanzable, y el frente.}
\Proc{\textup{\textsc{FrentePD}}$(A_1, \dots, A_n)$}{
  sea $M[0..2n][0..n][0..n]$ una tabla de costos; $M \gets \langle +\infty\rangle$;\quad $M[0][0][0] \gets 0$\;
  \For{$i \gets 1$ \KwTo $n$}{
    sea $M'$ una tabla como $M$; $M' \gets \langle +\infty\rangle$\;
    \ForEach{$\langle a, b, p\rangle$ con $M[a][b][p] < +\infty$}{
      \ForEach{$c \in A_i$}{
        $\langle a', b', p'\rangle \gets \langle a + \nu(c),\ b + \beta(c),\ p + \varphi_i(c)\rangle$\;
        \If{$M[a][b][p] + g_i(c) < M'[a'][b'][p']$}{
          $M'[a'][b'][p'] \gets M[a][b][p] + g_i(c)$\tcp*{se guarda $c$ para el retroceso}
        }
      }
    }
    $M \gets M'$\;
  }
  \Return{los no dominados de $\bigl\{\, \langle M[a][b][p], \tfrac{a}{2n}, \tfrac{b}{n}, \tfrac{p}{n}\rangle : M[a][b][p] < +\infty \,\bigr\}$}\;
}
\end{algorithm}

\fi
